\section{Projection and Retained Distinction}

\subsection{Declared projections}

\begin{definition}[Complete and projected state]
Let
\[
q:X\longrightarrow Y
\]
be a declared map between finite state spaces.  Relative to this map,
\(x\in X\) is the complete state and \(q(x)\in Y\) is its projected
state.
\end{definition}

The words complete and projected are map-relative.  They do not
assert that \(X\) is an absolute final ontology or that \(Y\) is the
only available registration.

\begin{definition}[Projection equivalence]
For \(x,x'\in X\), define
\[
x\sim_q x'
\quad\Longleftrightarrow\quad
q(x)=q(x').
\]
The equivalence classes are the fibres of \(q\).
\end{definition}

\begin{definition}[Projection-relative degeneracy]
A pair \(x\neq x'\) is projection-degenerate relative to \(q\) when
\[
q(x)=q(x').
\]
The complete distinction remains present in \(X\), while the declared
projection fails to resolve it.
\end{definition}

This is the exact finite meaning used in the present paper.  It is
not an assertion that two complete states have merged.

\subsection{Projected and complete return}

Let
\[
T:X\longrightarrow X
\]
be a declared finite transformation.

\begin{definition}[Return predicates]
Complete return at step \(m\) means
\[
T^{m}x=x.
\]
Projected return at step \(m\) means
\[
q(T^{m}x)=q(x).
\]
\end{definition}

\begin{proposition}[Registered-return implication]
For every declared finite transformation and projection,
\[
T^{m}x=x
\quad\Longrightarrow\quad
q(T^{m}x)=q(x).
\]
The converse need not hold.
\end{proposition}

\begin{proof}
The implication follows by applying \(q\) to the equality
\(T^{m}x=x\).  The failure of the converse is realized explicitly by
the certified lifted and registered recurrence examples imported from
the earlier theorem chain.
\end{proof}

Two established examples are:

\begin{enumerate}[label=(\roman*)]
\item a three-step projected triangle return whose twisted lift
requires six steps for complete return;

\item a ten-step quotient return in the registered five-frame
construction whose complete registered return requires twenty steps.
\end{enumerate}

These examples act on different finite objects and have different
receipt groups.  Only the structural implication is shared.

\subsection{Registration data and registered history}

\begin{definition}[Registration datum]
For a declared projection \(q:X\to Y\), a registration datum is
finite information retained in the complete construction that is not
determined by the projected state \(q(x)\).
\end{definition}

\begin{definition}[Registered history]
A registered history is an ordered finite composition of declared
transformations together with the retained registration data required
to distinguish the resulting complete state from its projected state.
\end{definition}

The established construction contains several examples:

\begin{itemize}
\item deck receipts in finite covers;
\item binary closure data;
\item binary edge gauge classes;
\item integer refinements;
\item registered frame coordinates;
\item primitive registered-history cycle classes.
\end{itemize}

These are not one universal receipt group.  Their common role is that
they preserve distinctions erased by a specified forgetful map.

\subsection{Scalar compression is another projection}

Let
\[
\rho:S\longrightarrow\mathbb R^n
\]
be a declared real-valued readout from a richer state space \(S\).
The composition
\[
S
\overset{\rho}{\longrightarrow}
\mathbb R^n
\overset{R_n}{\longrightarrow}
\mathbb R_{\ge0}
\]
is generally many-to-one twice over.

Distinct complete states may have the same vector readout.  Distinct
centered vectors may also have the same centered magnitude.  Thus
\[
R_n(\rho(s))=R_n(\rho(s'))
\]
does not imply
\[
s=s'.
\]

This provides an exact finite model for loss of registered
individuation without loss of complete distinction.

\subsection{No universal information ladder}

The finite constructions do not supply one canonical chain
\[
\text{complete state}
\longrightarrow
\text{lift}
\longrightarrow
\text{registration}
\longrightarrow
\text{scalar}
\]
covering every object in the theory.  They supply several declared
forgetful maps with separately certified domains.

Accordingly, every saturation claim in this paper must state:

\begin{enumerate}[label=(\roman*)]
\item the complete state space;
\item the declared projection or readout;
\item the distinction being retained;
\item the distinction being forgotten.
\end{enumerate}
